% Discussion（草稿）
\label{sec:discussion}

\subsection{Why the Market Remains Unbeaten}
Our results are consistent with semi-strong market efficiency in
football betting: closing odds aggregate information more accurately
than any single model we evaluated. The 87\% share of predictive
information carried by market features, the failure of value betting
even after calibration, and the absence of exploitable
opening-to-closing drift all point in the same direction. The
practical implication is not that forecasting models are useless,
but that their contribution must be measured in risk-adjusted terms
rather than raw accuracy.

The draw is the clearest example of a task-inherent limitation:
both our models and the betting market itself under-predict draws by
an order of magnitude (predicted draw rates of 2.3\% for XGBoost and
0.3\% for the market versus an empirical rate of 25.5\%). Draws are
the least frequent and least distinctive outcome class in football;
this low signal-to-noise ratio bounds the achievable macro-F1 of any
outcome classifier on this task, independently of model choice.

\subsection{Overconfidence, Not Information}
The value-betting experiment is instructive. Bets selected because
the model's expected value was positive lost 11.5\% ROI, and
restricting to larger positive EV made losses worse. After isotonic
calibration the losses shrank but remained negative. The model's
``edge'' is concentrated exactly where it is most confident and most
wrong---a signature of overconfidence rather than information. This
contrast---uncertainty filtering helps, confidence-based selection
hurts---is the central practical insight of the paper.

\subsection{Real-World Frictions}
All financial results are gross of transaction costs. Bookmaker
margins are already embedded in the odds we settle against, but
real-world execution would add betting taxes, withdrawal fees, and
odds slippage---typically 3--7\% of turnover depending on the
jurisdiction. The small positive ROIs we observe (up to 1.8\% for the
UI-tier strategy) would be entirely consumed by such frictions. The
practical value of the risk-control framework is therefore not
\emph{profitability} but \emph{risk stratification}: identifying, in
advance, the subset of matches where even a well-calibrated model is
unlikely to be right, so that capital and attention are directed
elsewhere. Claims of stable, deployable profitability would be
unsupported by this evidence.

\subsection{Uncertainty as a Risk Signal}
The uncertainty index and scenario complexity score both stratify
held-out matches monotonically, and the pattern replicates across
two different model families (XGBoost and the LLM). The multi-sample
consistency term contributes little (see
Section~\ref{sec:llm_analysis}); the effective components of the UI
are predictive confidence and market volatility, both cheap and
available at prediction time.

\subsection{LLM Results in Context}
The domain-enhanced LLM reaches market-level probability calibration
(log loss 0.967 vs 0.967, ECE 0.025 vs 0.026) with no training and no
retrieval, at a cost of \$0.55 for the entire test season. Two
caveats keep this claim honest. First, multi-sample consistency---a
mechanism proposed in earlier LLM-forecasting work---carries no
information in our setting: at both temperature 0.3 and 0.7 the
sampled probabilities agree almost perfectly and do not correlate
with correctness. Second, an open-weight 7B model matches or exceeds
the API model on a 200-match subset, so the result is not a property
of a specific proprietary model. We therefore frame the LLM
contribution as \emph{calibration without training} and a
cost--performance reference point, not as predictive superiority.

\subsection{Limitations}
\begin{itemize}
\item Text modality: we deliberately exclude news text. Publicly
available news corpora lack reliable pre-match timestamps at scale;
a rigorous text experiment would require timestamped, match-linked
articles. We sketch the constraints such an extension would need:
(i)~a strict pre-match cutoff, with every article's publication time
verified against the kickoff time and any post-kickoff or undated
article discarded; (ii)~protection against retrieval contamination, so
the LLM cannot reproduce outcomes it may have seen during pretraining
or via web access (no retrieval at inference, as in this paper);
(iii)~reporting coverage, since news availability is highly
heterogeneous across leagues and matches. Without these controls,
text features would reintroduce exactly the leakage this paper is
designed to eliminate.
\item The test season (2025/26) is incomplete; results may shift
slightly at season end. This limitation is actionable: every number in
this paper is produced by a script-driven pipeline
(\texttt{refresh\_all.cmd} reruns feature construction, all
\texttt{run\_*.py} experiments, and table/figure generation from the
updated CSV files in a single pass), so the complete result set and
every table can be regenerated without manual editing once the season
finishes. The public repository documents the exact refresh procedure.
\item Financial results are gross of betting taxes and use a single
bookmaker's closing odds; real-world execution may differ.
\item The LLM evaluation uses a proprietary API (DeepSeek-Chat);
we additionally report an open-weight model comparison
(see \S\ref{sec:llm_analysis}).
\end{itemize}
